% Generated draft: requires mathematical and source review.
\documentclass[11pt]{article}
\usepackage[margin=25mm]{geometry}
\usepackage{fontspec}
\setmainfont{Latin Modern Roman}
\usepackage{amsmath,amssymb,mathtools,unicode-math}
\setmathfont{Latin Modern Math}
\usepackage{enumitem,needspace,xurl,hyperref}
\setlength{\parindent}{0pt}
\setlength{\parskip}{5pt}
\setlength{\emergencystretch}{3em}
\newcommand{\N}{\mathbb{N}}
\newcommand{\Z}{\mathbb{Z}}
\newcommand{\Q}{\mathbb{Q}}
\newcommand{\R}{\mathbb{R}}
\newcommand{\C}{\mathbb{C}}
\newcommand{\F}{\mathbb{F}}
\newcommand{\A}{\mathbb{A}}
\newcommand{\PP}{\mathbb{P}}
\newcommand{\E}{\mathbb{E}}
\newcommand{\eps}{\varepsilon}
\newcommand{\dd}{\,\mathrm{d}}
\begin{document}
\section*{UnsolvedMath ID: 30000629}
\textbf{Source problem code:} catalogue entry
\subsection{Definitions and mathematical statement}
A convex body in $\R^n$ is a compact convex set with nonempty interior. For such a body $K$, its covariogram is
\[
g_K:\R^n\to[0,\infty),\qquad
g_K(x)=\operatorname{vol}_n(K\cap(K+x)).
\]
The question is whether, in the dimension under consideration, every pair of convex bodies $K,L\subset\R^n$ satisfies
\[
\left[\forall x\in\R^n,\ g_K(x)=g_L(x)\right]
\quad\Longrightarrow\quad
\left[\exists a\in\R^n:\ L=K+a\ \text{or}\ L=-K+a\right].
\]
The alternatives are precisely translation and translation followed by a point reflection. They are unavoidable ambiguities of the covariogram. Equality of data is pointwise in the same coordinates; arbitrary rotations are not included in the equivalence relation.

The record poses the question without a restriction on $n$, so the answer must be assessed dimension by dimension or as a universal assertion over dimensions. The supplied broad low-dimensional/higher-dimensional status summary is retained as dated source content, not as a proof that the displayed implication is true or open in each dimension.
\subsection{Short English statement}
Does the full overlap-volume function of a convex body determine it up to translation and reflection through a point? The competitor may be any convex body in the same dimension, and the answer may depend on that dimension.
\subsection{Sources}
\begin{itemize}
\item Supplied report: \emph{Konvexgeometrie} (2007), Oberwolfach Reports, DOI 10.4171/owr/2006/56, pp. 3321--3396; proposed year 2006. \url{https://doi.org/10.4171/owr/2006/56}.
\item Local UnsolvedMath registry, permanent ID 30000629, catalogue entry; created 25 August 2026, updated 16 September 2026. The corrected\_verified note says the uniqueness qualifications were retained while notation was normalized. Its supplied literature triage of 21 August 2026 labels the problem partially solved and states broadly that low dimensions are solved while the full higher-dimensional problem remains open. That summary gives no dimension-specific theorem or counterexample qualifications and requires independent checking; no additional external reference is supplied.
\end{itemize}
\end{document}
